% No preamble. Prefix fqc-. All equation references are internal.
\section{Continuity in the real coupling at a finite regulator}
\label{fqc-section}

This section concerns the actual homogeneous local macrostep and its
specified common-domain continuation. It supplies continuity needed
for a finite endpoint argument. It proves no coupling sign, endpoint
crossing, or identification with the full Wilson pushforward.

\subsection{Fixed data and statement}
\label{fqc-setup}

Fix one allowed finite macrostep, including its preceding classical
depth $M$, current and output tori, ratios $R,J$, common angle radius
$b$, carrier exponent $a_0$, actual retained chart, and all strict
common-domain budgets. Use the finite activity quotients by equality
of vacuum-normalized totals on the full product ball. Their norms
are infima of anchored family norms.

Let $I$ be a compact interval contained in $(0,g_0]$, where $g_0$
meets the common-domain and two-input comparison bounds. For example,
take the continuous threshold
\begin{equation}
 T(g)=\max\{T_0,A\sqrt{\log(1/g)}\}.
 \label{fqc-threshold}
\end{equation}
More generally, any positive continuous threshold satisfying those
bounds suffices. Only continuity and boundedness on $I$ will be used.
Denote the actual finite quotient map by
\begin{equation}
 F_g(x)=T_{\mathrm{cl}}x+\Delta+\mathcal Q_{\det}
                         +\mathcal L_g(Q_x).
 \label{fqc-map}
\end{equation}
Here $T_{\mathrm{cl}}$ is classical background transport, the next two
terms are the fixed pivot and Gaussian determinant corrections, and
$Q_x$ is the inherited fluctuation difference. The stationary
backgrounds, precision and chart are independent of $g$. The
coupling is inserted \emph{inside} the nonlinear chart as
$\Phi_V(g\xi)$.

\begin{theorem}[Finite real-coupling continuity]
\label{fqc-theorem}
On the common admissible quotient ball, for instance
$\lVert x\rVert\leq\overline\epsilon/4$, the map
$(g,x)\mapsto F_g(x)$ is jointly continuous for $0<g\leq g_0$.
Both activity topologies use the same radius $b$ and exponent $a_0$.
Consequently the actual extracted update
\begin{equation}
 (\gamma,x)\longmapsto
 \left(\gamma-\beta_{\mathrm{out}}
                 (F_{\gamma^{-1/2}}(x)),
       P_{\mathrm{out}}F_{\gamma^{-1/2}}(x)\right)
 \label{fqc-extracted-map}
\end{equation}
is continuous on its admissible positive weak-coupling domain.
The finite maps $\beta_{\mathrm{out}},P_{\mathrm{out}}$ are normalized
by the actual finite reference and do not depend on $g$.
\end{theorem}

The proof permits constants depending exponentially on the finite
volume. It asserts neither a volume-uniform modulus nor continuity
at $g=0$. In particular it does not require a continuous choice of
near-minimal activity representatives.

\subsection{Strict inherited-field margin}
\label{fqc-margin-section}

Fix one input class $x$ and choose one admissible representative
$E$, held fixed as $g$ varies; one may take
$\lVert E\rVert<\overline\epsilon/2$.
The proved actual joint-family construction gives
\begin{equation}
 f_g(w,\xi)=\Omega_R(a_p(w),g\xi).
 \label{fqc-actual-field}
\end{equation}
Its absolute rooted family bound supplies, under the strict budgets,
\begin{equation}
 \lVert a_p(w)\rVert_\infty\leq b/4,\qquad
 \lVert f_g(w,\xi)\rVert_\infty\leq b/4
 \label{fqc-margin}
\end{equation}
for all $\lVert w\rVert_\infty<b$, $g\in I$, and every actual
retained link/color window of threshold $T(g)$. Fixed norm
conversions have already been included in the common-domain
constants. The argument below only requires a fixed inner radius
$sb$ with $s<1$.

Write $E_{\mathrm{raw}}(A)=\sum_H E_H(A|_H)$, and let
$V_{\mathrm{in}}$ be the number of input $J$-cells. The anchored norm
implies
\begin{equation}
 \sum_H\lVert E_H\rVert_b
       \leq V_{\mathrm{in}}\lVert E\rVert_{a_0;b}.
 \label{fqc-extensive-bound}
\end{equation}
Directional one-variable Cauchy on the product ball gives
\begin{equation}
 |E_{\mathrm{raw}}(A)-E_{\mathrm{raw}}(B)|
 \leq
 \frac{V_{\mathrm{in}}\lVert E\rVert_{a_0;b}}
      {(1-s)b}\,\lVert A-B\rVert_\infty,
 \qquad \lVert A\rVert_\infty,\lVert B\rVert_\infty\leq sb.
 \label{fqc-input-cauchy}
\end{equation}
Indeed the segment between $A,B$ stays in the inner product ball;
each unit max-norm direction has a complex disk of radius at least
$(1-s)b$. Apply Cauchy on any smaller disk, integrate the derivative
along the segment, sum using \eqref{fqc-extensive-bound}, and pass
to the limiting disk radius. Thus an arbitrary input need only be
bounded and holomorphic on its original $b$ ball. No extension
past that ball is assumed.

\subsection{Fixing the real integration window}
\label{fqc-integral-section}

Let $D$ be the total finite number of real retained scalar
coordinates. Let $\Xi_1\subset\mathbb R^D$ be the product of the
actual centered unit retained link/color balls. Its dilation
$T(g)\Xi_1$ is exactly the actual window; its boundary is
Lebesgue-null. Let $A(w)$ be the retained precision, with the
fixed analytic square-root determinant branch. Changing real
variables $\xi=T(g)\zeta$ gives the exact Gaussian expectation
\begin{align}
 Z_g(w)
 ={}&\frac{\det A(w)^{1/2}}{(2\pi)^{D/2}}\,T(g)^D
       \int_{\Xi_1}
       \exp\bigl\{
          -\tfrac12 T(g)^2\zeta^T A(w)\zeta
          +V_{\mathrm{base}}(w,T(g)\zeta;g)
 \notag\\[-2pt]
 &\hspace{39mm}
          +E_{\mathrm{raw}}(f_g(w,T(g)\zeta))
          -E_{\mathrm{raw}}(a_p(w))
       \bigr\}\,d\zeta .
 \label{fqc-integral}
\end{align}
The background-only pivot and determinant terms may remain outside
this expression as in \eqref{fqc-map}. Coupling powers, compact
gauge-insertion constants and flat Haar constants are separately
tracked vacuum factors. No contour deformation or replacement
by a whole-fiber Euclidean cutoff occurs. The chart argument is
$\Phi_V(gT(g)\zeta)$.

The independent classical domains used in constructing the actual
chart and base interaction give the following bounds on this
fixed regulator. The precision and its determinant branch are
bounded for all $\lVert w\rVert_\infty<b$. The local analytic
coefficient maps are bounded and uniformly Lipschitz in their
fluctuation arguments on a slightly larger fixed chart domain
containing every $gT(g)\zeta$. Moreover
$V_{\mathrm{base}}(w,T(g)\zeta;g)$ is uniformly bounded and
uniformly continuous in $g\in I$, uniformly over the full open
$b$ ball and $\zeta\in\Xi_1$.

The last assertion follows directly from the exact factors as
well. The classical and gauge Taylor remainders are analytic
functions of $g\xi$, divided by $g^2$ after removing the constant,
linear and quadratic terms. The remaining Haar and Jacobian
factors are analytic in the same argument. On a compact positive
$g$ interval all these expressions have bounded parameter
derivatives on the fixed strict chart margin. Their removable
singularities at zero coupling are not needed. An anchored bound
on a classical support series bounds its extensive sum at this
fixed volume.

Since $a_p$ is independent of $g$, continuity of $gT(g)$ and these
analytic bounds imply
\begin{equation}
 \sup_{\substack{\lVert w\rVert_\infty<b\\\zeta\in\Xi_1}}
 \bigl\lVert f_g(w,T(g)\zeta)
       -f_{g_*}(w,T(g_*)\zeta)\bigr\rVert_\infty
       \longrightarrow0.
 \label{fqc-field-continuity}
\end{equation}
Use \eqref{fqc-input-cauchy} for the inherited exponent. The
Gaussian term has the same uniform continuity because $A(w)$
is bounded and $T$ is continuous. Every exponent in
\eqref{fqc-integral} has a common finite absolute bound; the
prefactor is uniformly bounded and $\Xi_1$ has finite measure.
Uniform domination therefore proves
\begin{equation}
 \lVert Z_g-Z_{g_*}\rVert_{H^\infty(b)}
       \longrightarrow0 .
 \label{fqc-z-continuity}
\end{equation}
This is a supremum over the entire open $b$ ball, not merely
convergence on compact subballs.

\subsection{A lower bound on the full ball and logarithms}
\label{fqc-log-section}

The strict Gaussian theorem supplies nonvanishing throughout
the $b$ ball and a connected logarithm family with a common
anchored bound $H_I$ for $g\in I$. If $V_{\mathrm{out}}$ is the
number of output $J$-cells, its normalized total satisfies
\begin{equation}
 \sup_{\lVert w\rVert_\infty<b}
  |\log Z_g(w)-\log Z_g(0)|
       \leq 2V_{\mathrm{out}}H_I.
 \label{fqc-log-bound}
\end{equation}
The bound is unaffected by raw logarithm constants.
By \eqref{fqc-z-continuity} at $w=0$ and nonvanishing, on a
compact positive interval around $g_*$,
$m_0=\inf_{g\in I}|Z_g(0)|>0$. Hence
\begin{equation}
 \inf_{\substack{g\in I\\\lVert w\rVert_\infty<b}}
       |Z_g(w)|
 \geq m_0\exp(-2V_{\mathrm{out}}H_I)=:m_I>0.
 \label{fqc-lower-bound}
\end{equation}
Nonvanishing alone would not yield such an infimum on the open
ball; the uniform logarithm estimate is essential.

When $\lVert Z_g-Z_{g_*}\rVert_b<m_I/2$, use the convergent
power-series logarithm
\begin{equation}
 \ell_g(w)=
 \log\left(1+\frac{Z_g(w)-Z_{g_*}(w)}{Z_{g_*}(w)}\right),
 \qquad
 \lVert\ell_g\rVert_b\leq
          2m_I^{-1}\lVert Z_g-Z_{g_*}\rVert_b.
 \label{fqc-log-ratio}
\end{equation}
If $l_g$ is any selected holomorphic logarithm of $Z_g$, then
$l_g-l_{g_*}-\ell_g$ is constant on the connected product ball,
because its exponential is one. Thus
\begin{equation}
 \bigl\lVert[l_g-l_g(0)]-[l_{g_*}-l_{g_*}(0)]\bigr\rVert_b
 \leq 4m_I^{-1}\lVert Z_g-Z_{g_*}\rVert_b
 \longrightarrow0.
 \label{fqc-normalized-log-continuity}
\end{equation}
No continuity of a raw branch convention has been assumed.

\subsection{From totals to the finite quotient norm}
\label{fqc-quotient-section}

Let $H_{\mathrm{all}}$ be the full connected output $J$-cell
torus. Assign any bounded holomorphic normalized total $h$ to
this one carrier and zero to all others. This is an admissible
finite activity representation and gives
\begin{equation}
 \lVert[h]\rVert_{\mathrm{quot}}
       \leq e^{a_0V_{\mathrm{out}}}\lVert h\rVert_b.
 \label{fqc-single-carrier}
\end{equation}
If $h$ has invariant total, this representative retains that
property; coefficientwise gauge invariance is unnecessary.
The exterior terms in \eqref{fqc-map} are independent of $g$.
Applying \eqref{fqc-single-carrier} to
\eqref{fqc-normalized-log-continuity} proves, for fixed $x$,
\begin{equation}
 \lVert F_g(x)-F_{g_*}(x)\rVert_{\mathrm{quot}}
       \longrightarrow0.
 \label{fqc-fixed-input-continuity}
\end{equation}
The volume cost in \eqref{fqc-single-carrier} is intentional.
It avoids an unsupported inference about the norm continuity
of every canonical cluster coefficient or relocalized array.

\subsection{Joint continuity and finite histories}
\label{fqc-joint-section}

The finite two-input estimate is uniform on $I$:
\begin{equation}
 \lVert F_g(x)-F_g(y)\rVert_{\mathrm{quot}}
 \leq
 \left[B_T+C_{\mathrm{quant}}
                  \sup_{h\in I}\frac{hT(h)}b\right]
       \lVert x-y\rVert_{\mathrm{quot}}.
 \label{fqc-input-lipschitz}
\end{equation}
It follows from admissible near-minimal representatives
$Y,D,X=Y+D$ and the stronger comparison budget
$a_Q\overline\epsilon\leq\eta_*/(128e)$. It does not require
a continuous representative selection. For $(g,x)\to(g_*,x_*)$,
split the difference as
\[
 F_g(x)-F_{g_*}(x_*)
   =F_g(x)-F_g(x_*)+F_g(x_*)-F_{g_*}(x_*).
\]
Equations \eqref{fqc-input-lipschitz} and
\eqref{fqc-fixed-input-continuity} make both terms tend to
zero. Boundedness of the actual finite $\beta_{\mathrm{out}}$
and $P_{\mathrm{out}}$, together with continuity of
$\gamma\mapsto\gamma^{-1/2}$, proves
Theorem~\ref{fqc-theorem}.

Every fixed finite composition of the extracted updates is
continuous while all intermediate inputs and couplings remain
in their common admissible domains. This follows by induction
over the finite stages. One can work near a trajectory lying
strictly within these domains, or use the finite-induction
envelope to ensure admissibility on a prescribed parameter
interval.

Near-minimal representatives implementing the history may
jump with the initial coupling. Exact finite-total descent
makes the quotient update and its beta increment independent
of those choices. Raw vacuum factors remain separately
tracked by their exact compensation identities; an endpoint
statement about $\gamma$ does not require a continuous raw
representative convention.

Continuity alone supplies no endpoint crossing. An
intermediate-value argument also needs values on opposite
sides of the desired endpoint and control keeping the
intervening histories weak and admissible. A bound on the
magnitude of each increment may keep a sufficiently high
finite initial $\gamma$ weak, but it does not tune a fixed
terminal coupling or establish an ultraviolet trajectory.
The present conclusion remains finite-volume, positive-coupling,
and confined to the specified local analytic map.
